%%
%% This is file `example_it.tex',
%% generated with the docstrip utility.
%%
%% The original source files were:
%%
%% ufrj-coppe.dtx  (with options: `exampleit')
%% 
%% This is one of the demonstration documents of the `ufrj' document class with
%% the `ufrj-coppe' unit style: the five main-language demos and the
%% side-by-side cover sheet.  It is generated from ufrj-coppe.dtx by docstrip;
%% edit the .dtx source, not this file.
%% 
%% Copyright (C) 2008-2026 Vicente Helano F. Batista, George O. Ainsworth Jr.,
%% Geraldo Xexéo and Eduardo Mangeli. Released under the GNU General Public
%% License version 3 (see the COPYING.txt file).
%% 
\documentclass[italian,dsc]{ufrj}
\usepackage{ufrj-coppe}
\addbibresource{exemplo.bib}
\begin{document}
  \title{Título do Trabalho}
  \subtitle{um exemplo de uso do CoppeTeX em italiano}
  \foreigntitle{Title of the Work}
  \foreignsubtitle{an example of use of CoppeTeX in Italian}
  \titlein{italian}{Titolo del Lavoro}
  \subtitlein{italian}{un esempio d'uso di CoppeTeX in italiano}
  \author{Nome dell'}{Autore}
  \advisor{Primo}{Relatore}{D.Sc.}{UFRJ}
  \examiner{Primo Esaminatore}{D.Sc.}{UFRJ}
  \examiner{Secondo Esaminatore}{D.Sc.}{UFRJ}
  \department{PESC}
  \date{05}{2026}
  \keyword{Multilingue}
  \keyword{Dimostrazione}
  %% palavras-chave do resumo em idioma estrangeiro (3.1.2.1.5, Anexo F)
  \foreignkeyword{Multilingual}
  \foreignkeyword{Demonstration}
  %% palavras-chave do resumo em português, a língua vernácula (3.1.2.1.4)
  \braziliankeyword{multilíngue}
  \braziliankeyword{demonstração}
  % Ordine degli elementi pre-testuali (Manuale UFRJ 3.1.2), verificato dalla
  % classe: \maketitle, \frontmatter, il riassunto in portoghese -- la lingua
  % vernacola, sempre il primo; qui brazilianabstract --, il riassunto in
  % inglese (foreignabstract), il riassunto in italiano (abstract), gli
  % elenchi e \tableofcontents, sempre l'ultimo.
  \maketitle
  \frontmatter
  \begin{brazilianabstract}
    Este é o resumo em português, a língua vernácula, que abre os resumos
    qualquer que seja o idioma do trabalho (Manual 3.1.2). Quem o produz é o
    ambiente \texttt{brazilianabstract}, quando o idioma principal não é nem
    o português nem o inglês.
  \end{brazilianabstract}
  \begin{foreignabstract}
    This is the English foreign abstract. Required in CoppeTeX whichever
    the main language is, for international readability.
  \end{foreignabstract}
  \begin{abstract}
    Questo è il riassunto in italiano, lingua principale della tesi, che
    viene dopo i riassunti in portoghese e in inglese; la copertina è
    conservata in portoghese per convenzione istituzionale dell'UFRJ. Lorem
    ipsum dolor sit amet, consectetur adipiscing elit.
  \end{abstract}
  % Elenchi pre-testuali: solo l'indice generale è obbligatorio (Manuale
  % UFRJ 3.1.2.1.6). Gli elenchi di figure, tabelle, abbreviazioni e simboli
  % sono facoltativi (3.1.2.2); quelli di quadros, programmi e algoritmi
  % compaiono quando il lavoro ne contiene (norma COPPE, §10). Questa
  % dimostrazione stampa solo l'indice.
  \tableofcontents
  \mainmatter
  \chapter{Introduzione}
  Questo capitolo dimostra il corpo redatto in italiano, mentre la
  copertina e la folha-de-rosto rimangono in portoghese (modello
  istituzionale UFRJ).
  \section{Il file di deposito: PDF/A}
  Il deposito all'UFRJ è esclusivamente digitale dalla risoluzione CEPG
  n.~246/2023, e il Manuale dell'UFRJ (2.2d) richiede il file finale in
  \textbf{PDF/A}, il formato di archiviazione a lungo termine. Il PDF/A è il PDF
  ristretto dalla norma ISO~19005 per la conservazione: tutti i font sono
  incorporati, i colori dichiarano un profilo (sRGB) che li rende indipendenti
  dal dispositivo, i metadati stanno in un blocco XMP letto dai repository, e
  nulla dipende dall'esterno --- niente contenuti esterni, JavaScript o
  cifratura. Così il file si apre allo stesso modo tra decenni, su qualsiasi
  computer.

  La classe usa il \textbf{PDF/A-2b}: la parte 2 ammette la trasparenza, che la
  parte 1 vieta e che riquadri colorati e grafici producono spesso, e il livello
  \emph{b} garantisce una riproduzione visiva fedele. Basta l'opzione
  \texttt{pdfa}: la classe carica \texttt{pdfx}, costruisce i metadati da
  titolo, autore e parole chiave, e incorpora il profilo colore. Attivatela
  presto, come fa \texttt{max-exemplo.tex}. Per verificare, usate il validatore
  veraPDF, profilo PDF/A-2b. L'italiano non è una lingua ammessa per una tesi
  della COPPE (risoluzione CEPG n.~302/2024, art.~57): questo documento
  dimostra soltanto il meccanismo.

  \section{Flussi di scrittura}
  pdf\LaTeX{} tiene al massimo sedici file aperti in scrittura, e ogni elenco,
  indice o glossario ne occupa uno. Per questo la classe carica
  \texttt{morewrites}. L'opzione \texttt{semmorewrites} lo disattiva e fa
  risparmiare circa mezzo secondo per passata, ma solo per un lavoro senza
  indice analitico, senza più indici e senza glossario.
\end{document}
%% 
%%
%% End of file `example_it.tex'.
